\documentclass[a4paper,11pt]{article}
\usepackage[utf8]{inputenc}
\usepackage[T1]{fontenc}
\usepackage[french]{babel}
\usepackage{lmodern}
\usepackage{amsmath,amssymb,amsthm}
\usepackage{mathrsfs}
\usepackage{enumitem}
\usepackage{multicol}

\setlist[enumerate]{label=\textbf{\textcolor{Couleur8}{\arabic*.}}, left=1em}

% Si vous voulez aussi modifier les listes enumerate imbriquées :
\setlist[enumerate,2]{label=\textcolor{Couleur8}{\alph*)}}
\setlist[enumerate,3]{label=\textcolor{Couleur8}{\roman*)}}
\setlist[enumerate,4]{label=\textcolor{Couleur8}{\Alph*)}}
\usepackage{graphicx}
\usepackage{xcolor}
\usepackage{tcolorbox}
\usepackage{geometry}
\usepackage{fancyhdr}
\usepackage{titlesec}
\usepackage{cancel}
\usepackage{ifthen}
\usepackage{tikz}
\usetikzlibrary{arrows.meta}
\usepackage{tkz-tab}
\usepackage{eurosym}

% OPTION PROF/ÉLÈVE
% Décommentez UNE des deux lignes suivantes :
\newboolean{prof}\setboolean{prof}{true}   % Version professeur (avec corrections des devoirs)
%\newboolean{prof}\setboolean{prof}{false}  % Version élève (sans corrections des devoirs)

\definecolor{Couleur1}{HTML}{4A5D7A} %Th
\definecolor{Couleur2}{HTML}{A5B5D6} %Prop
\definecolor{Couleur3}{HTML}{A8C68A} %Méthode
\definecolor{Couleur6}{HTML}{8A9CC1} %Def
\definecolor{Couleur7}{HTML}{E6B459} %Rq Voc
\definecolor{Couleur8}{HTML}{D36740} %Titres
\definecolor{correction}{HTML}{D36740} % Couleur pour les corrections
\definecolor{coopmathscorr}{HTML}{F15929} % Couleur corrections coopmaths

% Configuration des marges
\geometry{
	top=2cm,
	bottom=2cm,
	inner=1.5cm,
	outer=1.5cm,
}

% Police
\renewcommand{\familydefault}{\sfdefault}

% Compteurs
\newcounter{seancenum}
\newcounter{seancenumD}
\setcounter{seancenum}{1}
\setcounter{seancenumD}{1}

% Configuration des en-têtes et pieds de page
\pagestyle{fancy}
\fancyhf{}
\ifthenelse{\boolean{prof}}{
    \fancyhead[L]{\textcolor{Couleur8}{\textbf{Corrections Automatismes - Classe de seconde - Version Professeur}}}
}{
    \fancyhead[L]{\textcolor{Couleur8}{\textbf{Corrections Automatismes - Séance 15 - Classe de seconde}}}
}
\fancyhead[R]{\textcolor{Couleur8}{\textbf{2025-2026}}}
\fancyfoot[L]{\textcolor{Couleur8}{Mme Bertail et Mme Pinto}}
\fancyfoot[R]{\textcolor{Couleur8}{\thepage}}
\renewcommand{\headrulewidth}{1pt}
\renewcommand{\headrule}{\hbox to\headwidth{\color{Couleur8}\leaders\hrule height \headrulewidth\hfill}}

% Environnements personnalisés
\newtcolorbox{seance}[1]{
    colback=Couleur6!10,
    colframe=Couleur1!65,
    fonttitle=\bfseries\large,
    title=Correction Séance \theseancenum #1,
    left=5mm,
    right=5mm,
    top=3mm,
    bottom=3mm,
    arc=0mm,
    after=\stepcounter{seancenum}
}

\newtcolorbox{seanceD}[1]{
    colback=Couleur6!5,
    colframe=Couleur6!45,
    fonttitle=\bfseries\large,
    title=Correction Devoirs - Séance \theseancenumD #1,
    left=5mm,
    right=5mm,
    top=3mm,
    bottom=3mm,
    arc=0mm,
    after=\stepcounter{seancenumD}
}

\begin{document}

\setcounter{seancenum}{15}
\newpage
\begin{seance}{} %15

\begin{enumerate}
\item Déterminer les valeurs interdites revient à déterminer les valeurs qui annulent le dénominateur du quotient, puisque la division par $0$ n'existe pas.\\
Or $4x-6=0$ si et seulement si $x=\dfrac{3}{2}$.\\
Donc l'ensemble des valeurs interdites est $\left\{\dfrac{3}{2}\right\}$.\\
Pour tout $x\in \mathbb{R}\smallsetminus\left\{\dfrac{3}{2}\right\}$,\\
$\begin{aligned}
\dfrac{-5+9x}{4x-6}&=0\\
-5+9x&=0\\
\end{aligned}$\\
$\dfrac{5}{9}$ n'est pas une valeur interdite, donc l'ensemble des solutions de cette équation est ${\color{correction}\boldsymbol{\mathscr{S}=\left\{\dfrac{5}{9}\right\}}}$.

\item $f$ est une fonction affine. Elle s'annule en $x_0=\dfrac{3}{4}$.\\
Comme $-4<0$, $f(x)$ est négatif pour $x>\dfrac{3}{4}$ et positif pour $x<\dfrac{3}{4}$.\\

\begin{tikzpicture}[baseline, scale=0.75]
\tkzTabInit[lgt=5,deltacl=0.8,espcl=3.5]{ $x$ / 2, $f(x)$ / 2}{ $-\infty$, $\dfrac{3}{4}$, $+\infty$}
\tkzTabLine{  , +, z, -}
\end{tikzpicture}

\item On utilise la formule $\dfrac{a^n}{a^p}=a^{n-p}$ avec $a=9$, $n=7$ et $p=2$.\\
$\dfrac{9^{7}}{9^{2}}=9^{7- 2}={\color{correction}\boldsymbol{9^{5}}}$

\item $66\text{ min}= 60 \text{ min}+6\text{ min}=1\text{ h}+\dfrac{1}{10}\text{ h}=1{,}1\text{ h}$.\\
Ainsi, $66$ min correspond à ${\color{correction}\boldsymbol{1{,}1}}$ {\color{correction}heure}.

\item Après une augmentation de $10\,\%$, le nouveau prix est $P \times 1{,}1$.\\
Après une deuxième augmentation de $10\,\%$, le prix devient :\\
$(P \times 1{,}1) \times 1{,}1 = {\color{correction}\boldsymbol{P \times \left(1 + \dfrac{10}{100}\right)^2}}={\color{correction}\boldsymbol{P \times  \left(1 + \dfrac{1}{10}\right)^2}} $\\
La bonne réponse est la réponse {\bfseries \color{correction}C}.

\end{enumerate}

\end{seance}

\end{document}